\subsection{同态模上的表示}

再次设 \(\mathfrak{g}_1\) 和 \(\mathfrak{g}_2\) 是 K 上的两个李代数，\(\mathbf{M}_i\) 是一个 \(\mathfrak{g}_i\)-模（\(i = 1, 2\)）。令 \(\mathrm{U}_i\) 为 \(\mathfrak{g}_i\) 的包络代数，\(\sigma_i\) 为从 \(\mathfrak{g}_i\) 到 \(\mathrm{U}_i\) 的典范映射。于是 \(\mathbf{M}_i\) 是左 \(\mathrm{U}_i\)-模，因而 \(\mathcal{L}_{\mathrm{K}}(\mathrm{M}_1, \mathrm{M}_2)\) 具有典范的左 \((\mathrm{U}_1^0 \otimes \mathrm{U}_2)\)-模结构。现在 \(\mathrm{U}_1^0 \otimes_{\mathrm{K}} \mathrm{U}_2\) 是 \(\mathfrak{g}_1^0 \times \mathfrak{g}_2\) 的包络代数，映射

\[
(x _ {1}, x _ {2}) \mapsto \sigma_ {1} (x _ {1}) \otimes 1 + 1 \otimes \sigma_ {2} (x _ {2})
\]

是从 \(\mathfrak{g}_1^0\times \mathfrak{g}_2\) 到此包络代数的典范映射。因此，在 \((\mathfrak{g}_1^0\times \mathfrak{g}_2)\) 上存在 \(\mathbf{M} = \mathcal{L}_{\mathbb{K}}(\mathbf{M}_1,\mathbf{M}_2)\)-模结构，使得

\[
\begin{array}{r l} ((x _ {1}, x _ {2}) _ {\mathrm{M}}. u) \cdot m _ {1} & = ((\sigma_ {1} (x _ {1}) \otimes 1 + 1 \otimes \sigma_ {2} (x _ {2})) \cdot u) \cdot m _ {1} \\ & = u ((x _ {1}) _ {\mathrm{M} _ {1}}. m _ {1}) + (x _ {2}) _ {\mathrm{M} _ {2}}. u (m _ {1}) \end{array}\tag{4}
\]

对所有 \(u \in \mathcal{L}_{\mathbb{K}}(\mathrm{M}_1, \mathrm{M}_2)\)、\(m_1 \in \mathrm{M}_1\) 均成立。此结构定义了 \(\mathfrak{g}_1^0 \times \mathfrak{g}_2\) 在 \(\mathbf{M}\) 上的表示。

若现在 \(\mathfrak{g}_1 = \mathfrak{g}_2 = \mathfrak{g}\)，则从 \(x \mapsto (-x, x)\) 到 \(\mathfrak{g}\) 的同态 \(\mathfrak{g}^0 \times \mathfrak{g}\) 与上述表示复合，定义了 g 在 M 上的表示，因而在 M 上定义了 g-模结构，使得

\[
(x _ {\mathrm{M}}. u). m _ {1} = x _ {\mathrm{M} _ {2}}. u (m _ {1}) - u (x _ {\mathrm{M} _ {1}}. m _ {1})\tag{5}
\]

或

\[
x _ {\mathrm{M}}. u = x _ {\mathrm{M} _ {2}} u - u x _ {\mathrm{M} _ {1}}.\tag{6}
\]

结合这些结果与命题 2，得到：

命题 3. 设 \(\mathfrak{g}\) 是 \(\mathbf{K}\) 上的李代数，\(\mathbf{M}_i\) 是 \(\mathfrak{g}\)-模（\(1 \leqslant i \leqslant n + 1\)）。

令 \(\mathbf{N}\) 为 K-模 \(\mathcal{L}_{\mathbf{K}}(\mathbf{M}_1, \ldots, \mathbf{M}_n; \mathbf{M}_{n+1})\)，即从 \(\prod_{i=1}^{n} \mathbf{M}_i\) 到 \(\mathbf{M}_{n+1}\) 的多重线性映射的 K-模。在 \(\mathbf{N}\) 上存在唯一的 g-模结构，使得

\[
\begin{array}{r l} (x _ {N}. u) (m _ {1}, \ldots , m _ {n}) & = - \sum_ {i = 1} ^ {n} u (m _ {1}, \ldots , x _ {M _ {i}}. m _ {i}, \ldots , m _ {n}) \\ & \quad + x _ {M _ {n + 1}}. u (m _ {1}, \ldots , m _ {n}) \end{array}\tag{7}
\]

对所有 \(x \in \mathfrak{g}\)、\(u \in \mathbf{N}\) 及 \(m_i \in \mathbf{M}_i\)（\(1 \leqslant i \leqslant n\)）成立。

特别地，设 g 是 K 上的李代数，M 是 g-模，并把 K 看作平凡 g-模。命题 3 在 \(\mathcal{L}_{\mathrm{K}}(\mathrm{M},\mathrm{K}) = \mathrm{M}^{*}\) 上定义了 g-模结构。相应的表示称为表示 \(x\mapsto x_{\mathrm{M}}\) 的对偶表示。我们有：

\[
(x _ {\mathrm{M} ^ {*}}. f) (m) = - f (x _ {\mathrm{M}}. m)\tag{8}
\]

对所有 \(x \in \mathfrak{g}, f \in \mathbf{M}^*\)、\(m \in \mathbf{M}\) 成立。换言之：

\[
x _ {\mathrm{M} ^ {*}} = - ^ {t} x _ {\mathrm{M}}.\tag{9}
\]

当 K 是域且 M 是有限维时，g-模 M 是简单的（分别为半单的），当且仅当 g-模 M* 是简单的（分别为半单的）。

命题 4. 设 \(\mathbf{M}_1, \mathbf{M}_2\) 是两个 g-模。典范的 K-线性映射（《代数》，第 II 章，§ 4，第 2 号命题 2 以及第 1 号命题 1）：

\[
\mathrm{M} _ {1} ^ {*} \otimes_ {\mathrm{K}} \mathrm{M} _ {2} \xrightarrow {\Phi} \mathscr {L} _ {\mathrm{K}} (\mathrm{M} _ {1}, \mathrm{M} _ {2}), \quad \mathscr {L} _ {\mathrm{K}} (\mathrm{M} _ {1}, \mathrm{M} _ {2} ^ {*}) \xrightarrow {\Psi} (\mathrm{M} _ {1} \otimes_ {\mathrm{K}} \mathrm{M} _ {2}) ^ {*}
\]

（其中第二个是双射）都是 \(\mathfrak{g}\)-模同态。

  记作

\[
\mathrm{N} = \mathrm{M} _ {1} ^ {*} \otimes \mathrm{M} _ {2}, \quad \mathrm{P} = \mathscr {L} (\mathrm{M} _ {1}, \mathrm{M} _ {2}), \quad \mathrm{Q} = \mathscr {L} (\mathrm{M} _ {1}, \mathrm{M} _ {2} ^ {*}), \quad \mathrm{R} = (\mathrm{M} _ {1} \otimes \mathrm{M} _ {2}) ^ {*}.
\]

于是，对 \(x \in \mathfrak{g}, f \in \mathbf{M}_1^*\)、\(m_1 \in \mathbf{M}_1\)、\(m_2 \in \mathbf{M}_2\)，

\[
\begin{array}{r l} & {((\phi x _ {\mathrm{N}}) (f \otimes m _ {2})). m _ {1} = (\phi (x _ {\mathrm{M} _ {1} ^ {*}} f \otimes m _ {2} + f \otimes x _ {\mathrm{M} _ {2}} m _ {2})). m _ {1}} \\ & {\qquad = \langle x _ {\mathrm{M} _ {1} ^ {*}} f, m _ {1} \rangle m _ {2} + \langle f, m _ {1} \rangle x _ {\mathrm{M} _ {2}} m _ {2}} \\ & {((x _ {\mathrm{P}} \phi) (f \otimes m _ {2})). m _ {1} = x _ {\mathrm{M} _ {2}} (\phi (f \otimes m _ {2}). m _ {1}) - \phi (f \otimes m _ {2}) (x _ {\mathrm{M} _ {1}} m _ {1})} \\ & {\qquad = \langle f, m _ {1} \rangle x _ {\mathrm{M} _ {2}} m _ {2} - \langle f, x _ {\mathrm{M} _ {1}} m _ {1} \rangle m _ {2}} \end{array}
\]

从而 \(\phi x_{\mathbf{N}} = x_{\mathbb{P}}\phi\)。另一方面，对 \(x \in \mathfrak{g}\)、\(u \in \mathcal{L}(\mathrm{M}_1, \mathrm{M}_2^*)\)、\(m_1 \in \mathrm{M}_1\)、\(m_2 \in \mathrm{M}_2\)：

\[
\begin{array}{r l} (\psi x _ {\mathbf {Q}} u) (m _ {1} \otimes m _ {2}) & = \langle (x _ {\mathbf {Q}} u). m _ {1}, m _ {2} \rangle = \langle x _ {\mathrm{M} _ {2} ^ {*}} u m _ {1} - u x _ {\mathrm{M} _ {1}} m _ {1}, m _ {2} \rangle \\ (x _ {\mathrm{R}} \psi u) (m _ {1} \otimes m _ {2}) & = - \langle \psi u, x _ {\mathrm{M} _ {1}} m _ {1} \otimes m _ {2} + m _ {1} \otimes x _ {\mathrm{M} _ {2}} m _ {2} \rangle \\ & = - \langle u x _ {\mathrm{M} _ {1}} m _ {1}, m _ {2} \rangle - \langle u m _ {1}, x _ {\mathrm{M} _ {2}} m _ {2} \rangle \end{array}
\]

从而 \(\psi x_{Q} = x_{R}\psi\)，命题得证。

在同构 \(\mathcal{L}(\mathrm{M}_1,\mathrm{M}_2^*)\) 下，g-模 \((\mathrm{M}_1\otimes \mathrm{M}_2)^*\) 与 \(\psi\) 等同。若 \(\mathbf{M}_1\) 和 \(\mathbf{M}_2\) 有限基，则 \(\phi\) 是同构（《代数》，第 II 章，§ 4，第 2 号命题 2），这使我们可以把 g-模 \(\mathbf{M}_1^*\otimes \mathbf{M}_2\) 与 \(\mathcal{L}(\mathrm{M}_1,\mathrm{M}_2)\) 等同；在这种情况下，我们还可以把 g-模 \(\mathbf{M}_1^*\otimes \mathbf{M}_2^*\)、\(\mathcal{L}(\mathrm{M}_1,\mathbf{M}_2^*)\) 与 \((\mathrm{M}_1\otimes \mathrm{M}_2)^*\) 等同。

